\documentclass[letterpaper]{article} % DO NOT CHANGE THIS
\usepackage[submission]{aaai2027}  % DO NOT CHANGE THIS
% The serif, sans-serif, and monospaced fonts are loaded automatically by
% aaai2027.sty (newtxtext, helvet, courier). DO NOT add \usepackage{times},
% \usepackage{helvet}, \usepackage{courier}, or any other font package.
\usepackage[hyphens]{url}  % DO NOT CHANGE THIS
\usepackage{graphicx} % DO NOT CHANGE THIS
\urlstyle{rm} % DO NOT CHANGE THIS
\def\UrlFont{\rm}  % DO NOT CHANGE THIS
\usepackage{natbib}  % DO NOT CHANGE THIS AND DO NOT ADD ANY OPTIONS TO IT
\usepackage{caption} % DO NOT CHANGE THIS AND DO NOT ADD ANY OPTIONS TO IT
\frenchspacing  % DO NOT CHANGE THIS
%
% These are recommended to typeset algorithms but not required. See the subsubsection on algorithms. Remove them if you don't have algorithms in your paper.
\usepackage{algorithm}
\usepackage{algorithmic}

%
% These are recommended to typeset listings but not required. See the subsubsection on listing. Remove this block if you don't have listings in your paper.
\usepackage{newfloat}
\usepackage{listings}
\DeclareCaptionStyle{ruled}{labelfont=normalfont,labelsep=colon,strut=off} % DO NOT CHANGE THIS
\lstset{%
	basicstyle={\footnotesize\ttfamily},% footnotesize acceptable for monospace
	numbers=left,numberstyle=\footnotesize,xleftmargin=2em,% show line numbers, remove this entire line if you don't want the numbers.
	aboveskip=0pt,belowskip=0pt,%
	showstringspaces=false,tabsize=2,breaklines=true}
\floatstyle{ruled}
\newfloat{listing}{tb}{lst}{}
\floatname{listing}{Listing}

%
% Recommended for better-looking tables and mathematical notation
\usepackage{booktabs}
\usepackage{amsmath,amssymb}
\usepackage{multirow}
\usepackage{enumitem}

%
% Keep the \pdfinfo as shown here. There's no need
% for you to add the /Title and /Author tags.
\pdfinfo{
/TemplateVersion (2027.1)
}

% DISALLOWED PACKAGES
% \usepackage{authblk} -- This package is specifically forbidden
% \usepackage{balance} -- This package is specifically forbidden
% \usepackage{CJK} -- This package is specifically forbidden
% \usepackage{float} -- This package is specifically forbidden
% \usepackage{flushend} -- This package is specifically forbidden
% \usepackage{fullpage} -- This package is specifically forbidden
% \usepackage{geometry} -- This package is specifically forbidden
% \usepackage{hyperref} -- This package is specifically forbidden
% \usepackage{navigator} -- This package is specifically forbidden
% (or any other package that embeds links such as navigator or hyperref)
% \usepackage{indentfirst} -- This package is specifically forbidden
% \usepackage{layout} -- This package is specifically forbidden
% \usepackage{multicol} -- This package is specifically forbidden
% \usepackage{nameref} -- This package is specifically forbidden
% \usepackage{savetrees} -- This package is specifically forbidden
% \usepackage{setspace} -- This package is specifically forbidden
% \usepackage{stfloats} -- This package is specifically forbidden
% \usepackage{tabu} -- This package is specifically forbidden
% \usepackage{titlesec} -- This package is specifically forbidden
% \usepackage{tocbibind} -- This package is specifically forbidden
% \usepackage{ulem} -- This package is specifically forbidden
% \usepackage{wrapfig} -- This package is specifically forbidden
% DISALLOWED COMMANDS
% \nocopyright -- Your paper will not be published if you use this command
% \addtolength -- This command may not be used
% \balance -- This command may not be used
% \baselinestretch -- Your paper will not be published if you use this command
% \clearpage -- No page breaks of any kind may be used for the final version of your paper
% \columnsep -- This command may not be used
% \newpage -- No page breaks of any kind may be used for the final version of your paper
% \pagebreak -- No page breaks of any kind may be used for the final version of your paper
% \pagestyle -- This command may not be used
% \tiny -- This is not an acceptable font size.
% \vspace{- -- No negative value may be used in proximity of a caption, figure, table, section, subsection, subsubsection, or reference
% \vskip{- -- No negative value may be used to alter spacing above or below a caption, figure, table, section, subsection, subsubsection, or reference

\setcounter{secnumdepth}{0} %May be changed to 1 or 2 if section numbers are desired.

% The file aaai2027.sty is the style file for AAAI Press
% proceedings, working notes, and technical reports.
%

% Title

% Your title must be in mixed case, not sentence case.
% That means all verbs (including short verbs like be, is, using,and go),
% nouns, adverbs, adjectives should be capitalized, including both words in hyphenated terms, while
% articles, conjunctions, and prepositions are lower case unless they
% directly follow a colon or long dash
\title{DriveToken: Learning Dynamic Vision Token Budget with Driving Reward for Autonomous VLAs}
\author{
    Anonymous Submission
}
\affiliations{
}

%Example, Single Author, ->> remove \iffalse,\fi and place them surrounding AAAI title to use it
\iffalse
\title{My Publication Title --- Single Author}
\author {
    Author Name
}
\affiliations{
    Affiliation\\
    Affiliation Line 2\\
    name@example.com
}
\fi

\iffalse
%Example, Multiple Authors, ->> remove \iffalse,\fi and place them surrounding AAAI title to use it
\title{My Publication Title --- Multiple Authors}
\author {
    % Authors
    First Author Name\textsuperscript{\rm 1,\rm 2}\equalcontrib,
    Second Author Name\textsuperscript{\rm 2}\equalcontrib,
    Third Author Name\textsuperscript{\rm 1}\corresponding
}
\affiliations {
    % Affiliations
    \textsuperscript{\rm 1}Affiliation 1\\
    \textsuperscript{\rm 2}Affiliation 2\\
    firstAuthor@affiliation1.com, secondAuthor@affilation2.com, thirdAuthor@affiliation1.com
}
\fi


\begin{document}

\maketitle

\begin{abstract}
Vision-Language-Action (VLA) models for autonomous driving allocate substantial computation to vision tokens, although the visual evidence required for planning varies by scene. We present DriveToken, a two-stage plug-and-play framework that learns both which tokens to retain and how many tokens to retain. Stage~1 distills a lightweight 0.6M-parameter token scorer from internal attention using LambdaRank. Stage~2 jointly updates the token scorer and a scene-wise budget head with a REINFORCE objective that combines driving quality, safety-degradation penalties, and an efficiency incentive. The latest consolidated dynamic-budget configuration is evaluated on the complete NAVSIM navtest split (11,576 scenes) using AutoVLA-3B, physical token removal, a safety-net fallback, and a denylist patch. It reports PDMS \textbf{0.9107} (95\% CI $[0.9071,0.9143]$) at $35.5\%\pm12.0\%$ mean token retention ($256/720$ visual tokens), with 64.5\% visual-token and 43.3\% end-to-end FLOPs reduction. The fixed-ratio SFT scorer with safety-net fallback reports PDMS 0.9045 at $r{=}0.5$. Zero-shot transfer to ImpromptuVLA-7B yields pairwise ranking accuracy 0.856.
\end{abstract}


% Uncomment the following to link to your code, datasets, an extended version or similar.
% You must keep this block between (not within) the abstract and the main body of the paper.
% Make sure that you do not de-anonymize yourself with these links.
% \begin{links}
%     \link{Code}{https://aaai.org/example/code}
%     \link{Datasets}{https://aaai.org/example/datasets}
%     \link{Extended version}{https://aaai.org/example/extended-version}
% \end{links}
\newcommand{\ours}{DriveToken}
\newcommand{\best}[1]{\textbf{#1}}
\newcommand{\shard}{$^{\dagger}$}

\section{Introduction}
\label{sec:intro}

Vision-language-action (VLA) models have emerged as a promising paradigm for
end-to-end autonomous driving, unifying perception, reasoning, and action within
a single foundation model~\citep{openvla2024}.
A central bottleneck of this paradigm is computational: to reason over
multi-view camera inputs, a VLA encodes each view into a large grid of vision
tokens---e.g., 720 tokens from three surround cameras in AutoVLA---and the
subsequent LLM prefill over these tokens dominates inference cost (roughly 70\%
of total FLOPs).
Reducing the number of vision tokens is therefore one of the most direct levers
for making driving VLAs efficient and deployable.

Yet the importance of a vision token is fundamentally \emph{scene-dependent}.
In a safety-critical scenario---an occluded intersection with oncoming
traffic---tokens depicting the conflicting vehicle or the stop line are
indispensable, and aggressive pruning can be catastrophic.
In a simple highway-cruising scenario, the same spatial regions carry little
information and can be discarded with negligible loss.
The \emph{optimal token budget is not a fixed ratio but a property of the
driving context}, and a good pruning policy must decide \emph{both} which
tokens to keep \emph{and} how many to keep, per scene.

Prior visual-token reduction methods span three families, plus the emerging
driving-specific reconstruction approach.
\emph{Attention- and text-guided pruning}~\citep{chen2024fastv,xu2024sparsevlm}
selects tokens by cross-attention importance or text-token relevance at a
prescribed retention budget.
\emph{Similarity- and diversity-based pruning/merging}~\citep{shang2024llavaprumerge,topv2025,visionzip2025}
identifies redundant tokens through feature similarity or spatial diversity;
LLaVA-PruMerge~\citep{shang2024llavaprumerge} already adapts the retained
count to the input, yet its objective is visual-language coverage.
\emph{Foreground-reconstruction pruning for driving}~\citep{fastdrivevla2026}
targets background redundancy via reconstruction loss but deploys a fixed
Top-$K$ budget with no driving-reward signal.
\emph{Learned input-adaptive compression}~\citep{atpllava2025,glimpseprune2025}
varies the retained set per input via a trainable predictor, yet optimizes
grounding or VQA rather than a trajectory-quality driving objective.
The shared limitation is not that these methods are non-adaptive---several
are---but that \emph{none optimizes a trajectory-quality or safety objective}.
Token selection driven by attention, visual similarity, or reconstruction loss
can retain visually salient but driving-irrelevant tokens (sky, background
texture) and discard safety-critical ones (distant pedestrians, occluded
cyclists).

We propose \ours{}, a two-stage plug-and-play framework that grounds both
decisions---\emph{which} tokens to keep and \emph{how many}---in a
trajectory-quality driving objective, the first framework to our knowledge to do so.
\ours{} attaches a tiny 0.6M-parameter scorer and budget head to a frozen VLA
backbone with sub-millisecond overhead and no architectural change.
\emph{Stage\,1} trains the token scorer via \emph{LambdaRank distillation}
from the VLA's layer-12 attention: the listwise ranking objective teaches the
scorer to rank tokens by driving relevance, and because the scorer additionally
ingests camera position encodings it learns a ranking that \emph{surpasses its
own teacher} (PDMS 0.9045 vs.\ 0.8901 at matched retention).
\emph{Stage\,2} casts token budgeting as a reinforcement-learning problem: a
REINFORCE policy gradient optimizes a continuous budget head with a
three-component reward---(\textit{i})~shaped driving quality combining absolute
PDMS sub-metric scores and per-scene pruning deltas relative to the unpruned
baseline, (\textit{ii})~a safety-degradation penalty, and
(\textit{iii})~an efficiency bonus $\lambda_{e}(1{-}k_s)$ proportional
to token reduction.
The per-scene adaptivity of the budget follows from the joint driving-quality,
safety, and efficiency objective. At deployment, the evaluated dynamic
configuration additionally uses a safety-net fallback and a denylist patch.

The latest consolidated dynamic result is evaluated on all 11,576 NAVSIM navtest
scenes with AutoVLA-3B, physical token removal, five epochs of training, and
validation-based checkpoint selection. It reports PDMS
\textbf{0.9107} at an average 35.5\% token retention (256 of 720 tokens), with
43.3\% end-to-end FLOPs reduction. Its observed keep ratios span $[0.20,0.62]$;
59.5\% of scenes retain fewer than 40\% of visual tokens. The fixed-ratio SFT
scorer with safety-net fallback reports PDMS 0.9045 at $r{=}0.5$. \ours{} also
includes separate zero-shot ranking and cross-dataset analyses on 7B VLAs;
those experiments are reported independently from the NAVSIM dynamic budget
configuration.

\noindent Our contributions are:
\begin{itemize}[leftmargin=*,nosep]
\item \textbf{Driving-reward-guided joint token selection and budgeting.}
  \ours{} optimizes both the vision-token ranking and the per-scene retention
  budget under a trajectory-quality objective. The reward combines absolute PDMS
  sub-metric scores, same-scene pruning deltas, a safety-degradation penalty,
  and an efficiency term $\lambda_{e}(1{-}k_s)$. The latest safeguarded
  dynamic configuration reports PDMS \textbf{0.9107} at 35.5\% mean retention.

\item \textbf{Plug-and-play, zero-modification design.}
  The 0.6M MLP scorer and budget head attach to a frozen VLA backbone without
  altering weights or architecture, adding $<1$\,ms overhead. Physical token
  removal propagates efficiency savings through the full LLM prefill pass.

\item \textbf{Cross-model generalization and scaling law.}
  The scorer transfers zero-shot from AutoVLA-3B to ImpromptuVLA-7B (pairwise
  accuracy 0.856). Larger VLAs exhibit greater attention concentration,
  making \ours{}'s efficiency gains model-size scalable.
\end{itemize}
% 放在 \section{Related Work} 之后、第一个 \subsection 之前


\begin{figure}[t]
    \centering
    \includegraphics[width=\columnwidth]{Figures/1.png}
    \caption{\textbf{Overview of \ours{}.}
    Stage~1 learns which vision tokens to retain via LambdaRank-distilled
    driving-aware scoring. Stage~2 learns a scene-adaptive keep ratio through
    a Gaussian budget policy optimized with a three-component driving reward.
    Selected tokens are physically removed before frozen LLM decoding.}
    \label{fig:overview}
\end{figure}

\section{Related Work}
\label{sec:related}

\subsection{Visual Token Reduction for Vision-Language Models}

\textbf{Attention- and text-guided pruning.}
FastV~\citep{chen2024fastv} observes that visual-token attention decays sharply
after the first few LLM layers and removes low-attention tokens without
retraining.
SparseVLM~\citep{xu2024sparsevlm} uses text-token relevance to guide sparsity.
Both are training-free and efficient, but their pruning signal is visual
saliency or language relevance rather than driving quality.

\textbf{Similarity- and diversity-based pruning and merging.}
Token Merging (ToMe)~\citep{bolya2023tome} merges visually redundant tokens
based on feature similarity with minimal architectural change.
LLaVA-PruMerge~\citep{shang2024llavaprumerge} is \emph{adaptive}: the number
of retained tokens is determined by the IQR of CLS-attention outliers, followed
by similarity-based merging---its retained count varies per input.
VisionZip~\citep{visionzip2025} and TopV~\citep{topv2025} pursue coverage- and
diversity-aware selection via feature similarity or spatial structure, without
requiring a fixed budget.
These methods demonstrate that adaptive or diversity-driven compression can
preserve downstream accuracy, but their optimization objectives are
visual-language quality, not trajectory-quality driving performance.

\textbf{Learned input-adaptive compression.}
ATP-LLaVA~\citep{atpllava2025} and GlimpsePrune~\citep{glimpseprune2025} train
importance predictors that vary the retained token set per input.
Dynamic-LLaVA~\citep{dynamicllava2025} sparsifies context on the fly during
decoding.
These methods achieve per-input adaptivity but are designed and evaluated on
visual grounding and VQA, not trajectory-quality driving evaluation.

The gap shared across all three families is the \emph{objective}: none uses a
trajectory-quality or safety signal to learn either token ranking or budget.
\ours{} differs precisely here---the token scorer is distilled from a
driving-correlated teacher, and the budget head is trained with a PDMS reward.

% 放在 \section{Related Work} 之后、第一个 \subsection 之前
\begin{figure*}[t]
    \centering
    \includegraphics[width=\textwidth]{Figures/4.png}
    \caption{%
        \textbf{Taxonomy of prior visual token reduction methods.}
        \textbf{(a)}~\emph{Importance-/Salience-Driven}: tokens are scored
        by cross-attention, text relevance, or learned salience and the
        top-$K$ are retained at a fixed budget (FastV, SparseVLM,
        FasterVLM); FastDriveVLA learns salience via reconstruction but
        still deploys a fixed Top-$K$.
        \textbf{(b)}~\emph{Similarity-/Diversity-Driven}: redundant tokens
        are pruned or merged by feature similarity or spatial diversity
        (LLaVA-PruMerge, VisionZip, DivPrune); PruMerge uses IQR outliers
        making its budget input-adaptive, but optimizes visual-language
        coverage.
        \textbf{(c)}~\emph{Learned Input-Adaptive Compression}: a trainable
        predictor dynamically selects tokens per input (ATP-LLaVA,
        Dynamic-LLaVA, GlimpsePrune), yet optimizes VQA or generation
        quality rather than driving planning.
        \textbf{\ours{} (ours)} differs from all three: both the token
        ranking (\S\ref{sec:sft}) and the per-scene budget
        (\S\ref{sec:rl}) are optimized under a NAVSIM trajectory-quality
        objective.%
    }
    \label{fig:related_taxonomy}
\end{figure*}



\subsection{Token Reduction for Autonomous Driving}

\textbf{Driving VLMs.}
Prune2Drive~\citep{prune2drive2026} reduces cross-view redundancy in multi-camera
driving VLMs via token-wise farthest-point sampling; its retention budget is
allocated \emph{offline} per camera view and is not adapted per scene at
inference.
Its evaluation targets driving VLM reasoning rather than end-to-end trajectory
control.

\textbf{Driving VLAs.}
FastDriveVLA~\citep{fastdrivevla2026} introduces foreground-oriented token
selection via a reconstruction objective for driving VLA stacks, showing that
discarding background tokens can preserve or improve navigation scores.
However, it deploys a \emph{fixed} Top-$K$ budget at inference and does not
optimize a planning reward or learn a per-scene budget.

\textbf{Embodied VLAs (non-driving).}
VLA-Pruner~\citep{vlapruner2025} studies spatial-temporal dual-level token
selection for robotic manipulation VLAs.
It targets embodied manipulation rather than autonomous-driving control and does
not optimize a driving-quality metric.

\ours{} is the first, to our knowledge, to learn a continuous per-scene token
budget from a driving-quality reward in an end-to-end driving VLA, covering both
the \emph{what} (LambdaRank scorer) and the \emph{how many} (RL budget head)
under a planning objective.

\subsection{Reinforcement Learning for Efficient Inference}

Policy-gradient methods have been applied to discrete resource allocation,
e.g., early-exit depth selection~\citep{elbayad2020depth} and dynamic layer
execution~\citep{wang2018skipnet}.
GRPO~\citep{shao2024deepseekmath} introduces group-normalized advantage
estimation that stabilizes on-policy updates with small rollout batches.
\ours{} adapts GRPO to a \emph{continuous} budget action space via a Gaussian
policy over $[0.2,\,0.9]$, with an $L_2$ trust-region anchor to the SFT
initialization to prevent reward hacking on the token scorer.


\section{Method}
\label{sec:method}

\subsection{Problem Formulation and Overview}

We consider a frozen vision-language-action (VLA) policy $\mathcal{M}$ that
maps multi-view observations and a driving instruction $q$ to a future
trajectory. For each scene $s$, its vision encoder produces $N{=}720$ visual
tokens, arranged as $C{=}3$ surround-view cameras with 240 tokens per view.
Passing all tokens through the LLM decoder is costly; however, applying a fixed
retention ratio ignores the substantial variation in scene difficulty. Our
objective is therefore to learn a scene-conditional pruning action
\begin{equation}
    a_s = (k_s,\mathcal{S}_s), \qquad
    k_s \in [0.2,0.9], \quad
    |\mathcal{S}_s| = K_s = \operatorname{round}(k_sN),
\end{equation}
where $k_s$ specifies \emph{how many} tokens to retain and $\mathcal{S}_s$
specifies \emph{which} tokens to retain. The selected tokens are supplied to
the otherwise unchanged decoder.

\ours{} separates this decision into two stages. Stage~1 learns a lightweight
driving-aware token scorer from the frozen VLA's internal attention. Stage~2
treats the budget and the selected token set as a joint stochastic pruning
action and refines both using trajectory-quality feedback.
Figures~\ref{fig:stage1_scorer} and~\ref{fig:stage2_budget} detail the two
stages. The VLA backbone remains frozen throughout; only the small scorer and
budget head are optimized.

\subsection{Stage 1: Driving-Aware Token Scoring}
\label{sec:sft}

\paragraph{Scorer architecture.}
For token $i$, we use the layer-0 ViT-to-LLM projection
$\mathbf{v}_i\in\mathbb{R}^{2048}$ and its camera identity
$\mathbf{c}_i\in\{0,1\}^{3}$. The scorer input is
\begin{equation}
    \mathbf{x}_i =
    [\operatorname{Normalize}(\mathbf{v}_i);\mathbf{c}_i]
    \in\mathbb{R}^{2051}.
\end{equation}
Our scorer $f_\theta:\mathbb{R}^{2051}\to\mathbb{R}$ is a compact three-layer
MLP with approximately $0.6$M parameters:
\begin{equation}
 f_\theta(\mathbf{x}_i) =
 W_3\,\operatorname{GELU}\!\left(
 W_2\,\operatorname{GELU}\!\left(
 W_1\operatorname{LN}(\mathbf{x}_i)
 \right)\right),
\end{equation}
where $W_1\in\mathbb{R}^{256\times2051}$,
$W_2\in\mathbb{R}^{256\times256}$, and
$W_3\in\mathbb{R}^{1\times256}$. The camera encoding allows the scorer to
exploit view-specific driving priors while retaining a backbone-agnostic
interface. Scoring all 720 tokens requires a single lightweight forward pass
and adds negligible computation relative to VLA inference.

\paragraph{Attention teacher and ranking distillation.}
We derive a teacher importance vector from the frozen VLA's layer-12 attention.
Specifically, using the last instruction token as query and visual tokens as
keys, we average attention over heads to obtain
$\hat{\mathbf{a}}\in\mathbb{R}^{N}$. Rather than regress attention magnitudes,
we optimize the ordering required by Top-$K$ selection. For each teacher-ordered
pair $i\succ j$, LambdaRank minimizes
\begin{equation}
\mathcal{L}_{\mathrm{rank}} =
\sum_{i\succ j}\left|\Delta\operatorname{NDCG}_{ij}\right|
\log\!\left(
1+\exp\left[
-\big(f_\theta(\mathbf{x}_i)-f_\theta(\mathbf{x}_j)\big)
\right]\right),
\label{eq:lambdarank}
\end{equation}
where $\Delta\operatorname{NDCG}_{ij}$ is the change in NDCG induced by
swapping $i$ and $j$. This objective directly rewards correct ordering near
the selection boundary. We train this stage on 4,000 \textit{navtrain} scenes
and use the resulting scorer as the reference policy for Stage~2.

\begin{figure*}[t]
    \centering
    \includegraphics[width=\textwidth]
    {Figures/2.png}
    \caption{\textbf{Stage~1: driving-aware token scoring.}
    Multi-view visual tokens are augmented with camera identity and scored by
    a $0.6$M MLP. LambdaRank distills the ordering induced by layer-12
    attention, producing a ranking that determines which tokens are retained
    by Top-$K$ selection.}
    \label{fig:stage1_scorer}
\end{figure*}

\subsection{Stage 2: Safe Dynamic Budget Learning}
\label{sec:rl}

Stage~1 provides a strong ranking initialization, but attention-derived
relevance is only a proxy for driving utility. Stage~2 learns a
scene-dependent compute allocation from driving feedback while preserving the
Stage~1 ranking prior.

\paragraph{Stochastic budget policy.}
The budget head $g_\phi$ first mean-pools the token features,
$\bar{\mathbf{x}}_s=N^{-1}\sum_{i=1}^{N}\mathbf{x}_i$, and maps the scene
representation to a raw budget logit $\mu_\phi(\bar{\mathbf{x}}_s)$. During
training, we sample a latent action from a Gaussian policy,
\begin{equation}
 u_s \sim \mathcal{N}\!\left(
 \mu_\phi(\bar{\mathbf{x}}_s),\sigma_\phi^2\right),
 \qquad \sigma_\phi=\exp(\log\sigma_0),
\end{equation}
and squash it into the valid retention range:
\begin{equation}
 k_s = 0.2 + 0.7\,\operatorname{sigmoid}(u_s),
 \qquad K_s=\operatorname{round}(k_sN).
\label{eq:budget-policy}
\end{equation}
The stochastic policy encourages budget exploration; it does not impose a
fixed per-scene ratio. At inference, we deploy the policy mean in
Eq.~\ref{eq:budget-policy}, yielding deterministic but scene-dependent budgets
without sampling noise.

\paragraph{Joint selection--budget policy.}
Given $K_s$, the scorer chooses
\begin{equation}
\mathcal{S}_s =
\operatorname{TopK}
\left(\{f_\theta(\mathbf{x}_i)\}_{i=1}^{N},K_s\right).
\end{equation}
We use a multi-label softmax surrogate for the selected set's log-probability:
\begin{equation}
\log\widetilde{\pi}_\theta(\mathcal{S}_s\mid\mathbf{x}_s,k_s)
=
\frac{1}{K_s}\left[
\sum_{i\in\mathcal{S}_s} f_\theta(\mathbf{x}_i)
-K_s\log\sum_{j=1}^{N}
\exp\left(f_\theta(\mathbf{x}_j)\right)
\right].
\label{eq:selection-policy}
\end{equation}
The joint policy score is
\begin{equation}
\log\widetilde{\pi}_{\theta,\phi}(a_s\mid\mathbf{x}_s)
=
\log\pi_\phi(k_s\mid\bar{\mathbf{x}}_s)
+\lambda_{\mathrm{sel}}
\log\widetilde{\pi}_\theta
(\mathcal{S}_s\mid\mathbf{x}_s,k_s),
\label{eq:joint-policy}
\end{equation}
where $\lambda_{\mathrm{sel}}{=}1$ in our default configuration.
Equation~\ref{eq:selection-policy} is a low-variance Top-$K$ policy surrogate:
it provides a score-function gradient to the token scorer while retaining
deterministic Top-$K$ execution. Thus, the same trajectory-level feedback
updates both token ranking and scene-level budget allocation.

\paragraph{Driving-quality, efficiency, and safety reward.}
For a pruned rollout, let $q_m^{\mathrm{pruned}}$ and
$q_m^{\mathrm{base}}$ be the six NAVSIM submetrics for metric
$m\in\mathcal{M}$, where $\mathcal{M}$ contains collision avoidance,
drivable-area compliance, progress, time-to-collision (TTC), direction
compliance, and comfort. We first define a shaped driving reward,
\begin{equation}
R_{\mathrm{drive}}=
\alpha\sum_{m\in\mathcal{M}}w_mq_m^{\mathrm{pruned}}
+\beta\sum_{m\in\mathcal{M}}w_m
\left(q_m^{\mathrm{pruned}}-q_m^{\mathrm{base}}\right),
\label{eq:driving-reward}
\end{equation}
with $(\alpha,\beta)=(0.3,0.7)$ and
\begin{equation}
(w_{\mathrm{progress}},w_{\mathrm{collision}},
w_{\mathrm{drivable}},w_{\mathrm{TTC}},
w_{\mathrm{direction}},w_{\mathrm{comfort}})
=
(0.35,0.20,0.15,0.15,0.10,0.05).
\end{equation}
The absolute term accounts for scene difficulty, whereas the delta term
isolates pruning-induced degradation or improvement.

To discourage unsafe over-pruning, we explicitly penalize harmful degradation
on safety-critical submetrics:
\begin{equation}
L_{\mathrm{safe}} =
\sum_{m\in\{\mathrm{collision},\mathrm{drivable},\mathrm{TTC}\}}
\omega_m
\left[
q_m^{\mathrm{base}}-q_m^{\mathrm{pruned}}-\delta_m
\right]_+,
\label{eq:safety-loss}
\end{equation}
where
$(\omega_{\mathrm{collision}},\omega_{\mathrm{drivable}},
\omega_{\mathrm{TTC}})=(0.4,0.3,0.3)$ and
$[z]_+=\max(z,0)$. The rollout reward is
\begin{equation}
R_s =
\lambda_dR_{\mathrm{drive}}
-\lambda_sL_{\mathrm{safe}}
+\lambda_e(1-k_s),
\label{eq:total-reward}
\end{equation}
where the final term rewards compute reduction without prescribing a particular
$k_s$. In our pre-registered AAAI configuration,
$(\lambda_d,\lambda_s,\lambda_e)=(2.0,0.5,0.15)$.

\paragraph{Group-normalized policy optimization.}
For a group of $G$ rollout scenes, we form normalized advantages
\begin{equation}
\widehat{A}_s =
\frac{R_s-\operatorname{mean}(R_{1:G})}
{\operatorname{std}(R_{1:G})+\epsilon}.
\end{equation}
The policy objective is
\begin{equation}
\mathcal{L}_{\mathrm{PG}}
=
-\frac{1}{G}\sum_{s=1}^{G}
\widehat{A}_s
\log\widetilde{\pi}_{\theta,\phi}(a_s\mid\mathbf{x}_s).
\label{eq:pg}
\end{equation}
We anchor the token scorer, but not the budget head, to the Stage~1
initialization with a weight-space trust-region penalty,
\begin{equation}
\mathcal{L}_{\mathrm{ref}} =
\beta_{\mathrm{ref}}
\left\|\theta-\theta_{\mathrm{SFT}}\right\|_2^2,
\qquad \beta_{\mathrm{ref}}=0.01,
\end{equation}
and optimize
$\mathcal{L}=\mathcal{L}_{\mathrm{PG}}+\mathcal{L}_{\mathrm{ref}}$.
The frozen VLA is never back-propagated through. We synchronize gradients of
the scorer, budget head, and policy standard deviation across eight workers, so
all workers optimize a single shared policy.

\begin{figure*}[t]
    \centering
    \includegraphics[width=\textwidth]
    {Figures/3.png}
    \caption{\textbf{Stage~2: safe dynamic budget learning.}
    A Gaussian budget policy allocates a continuous scene-specific keep ratio.
    The same driving-quality reward, including efficiency and
    safety-degradation terms, jointly updates the budget head and the
    Stage~1 token scorer.}
    \label{fig:stage2_budget}
\end{figure*}

\subsection{Physical Token Removal for Evaluation and Deployment}
\label{sec:token-drop}

We use attention masking only as a fixed-shape rollout proxy during training.
All reported NAVSIM fixed-ratio and dynamic evaluations use physical token
removal: unselected visual-token entries are removed from the decoder inputs,
position IDs, attention mask, and cache positions before generation.
Consequently, LLM prefill and KV-cache cost scale with the retained
scene-specific $K_s$, rather than with the original 720 visual tokens.

\paragraph{Confidence-based runtime fallback.}
To guard against overly aggressive pruning on difficult scenes, we apply a
confidence-based runtime fallback. The token scorer's output logits are used to
compute an entropy-based confidence score per scene; if the confidence falls
below a calibrated threshold, the scene is re-evaluated with the full 720-token
no-pruning baseline (denylist patch). In practice, this fallback is triggered on
only a small fraction of scenes ($\sim$1--2\%), providing a lightweight safety
net without substantially increasing average compute. The denylist further
hard-codes a set of known failure modes identified during development. Both
mechanisms are reported explicitly as part of the deployment protocol in the
experiments.



\section{Experiments}
\label{sec:exp}

\subsection{Setup}
\label{sec:setup}

\paragraph{Dataset.}
We train the scorer on the \emph{navtrain} split of \textbf{NAVSIM}~\citep{navsim2024}
and evaluate physical token removal on \emph{navtest} ($N{=}11{,}576$ scenes).
For cross-dataset transfer (\S\ref{sec:gen}) we additionally evaluate on
\textbf{nuScenes} open-loop planning using ImpromptuVLA-7B~\citep{impromptuvla2025}.

\paragraph{Metrics.}
The primary metric is \textbf{PDMS} ($0$--$1$), decomposed into No-at-fault
Collision (NC), Drivable Area Compliance (DAC), Progress (EP),
Time-to-Collision (TTC), Direction, and Comfort. We report retained visual-token
count and end-to-end FLOPs reduction. Relative PDMS is normalized to the
unpruned AutoVLA-3B reference of 0.89879.

\paragraph{Baselines.}
We compare against \textbf{No Prune} (720 tokens), \textbf{Random},
\textbf{FastV}~\citep{fastv2024}, \textbf{PruMerge}~\citep{shang2024llavaprumerge},
and \textbf{SparseVLM}~\citep{sparsevlm2025}. All NAVSIM rows use physical token
removal. Baseline rows report raw scores without a denylist; our rows use an
explicit safety-net fallback and a denylist patch.

\paragraph{Training.}
\textbf{Stage~1 (scorer).} We train the LambdaRank scorer on the
\emph{navtrain} split using AdamW (learning rate $3\times10^{-4}$, weight
decay $10^{-4}$), 20 maximum epochs, 64 scenes per batch, and 1{,}024 sampled
token pairs per scene. We use cosine learning-rate decay to $10^{-5}$ and an
80/20 train/validation split of \emph{navtrain}, with early stopping (patience
three) on validation pairwise accuracy. The selected scorer initializes Stage~2.

\textbf{Stage~2 (joint budget RL).} We jointly optimize the scorer and budget
head for five epochs with AdamW, an initial learning rate of $10^{-5}$, cosine
decay, and 100 warmup steps. The batch size is 32. The token scorer is
regularized with $\beta_{\mathrm{ref}}=0.01$, whereas the budget head has no
trust-region constraint. The reward combines the absolute six-submetric term, the per-scene
pruning-delta term, and the efficiency bonus in Eq.~\ref{eq:total-reward}. We
select the epoch-4, step-2{,}180 checkpoint by PDMS on a 500-scene held-out
validation subset of \emph{navtrain}.

\paragraph{Backbone \& implementation.}
All NAVSIM results use \textbf{AutoVLA-3B}; cross-model generalization
additionally uses ImpromptuVLA-7B. Both fixed-ratio and dynamic evaluations use
physical token removal. The dynamic BudgetRL configuration additionally uses a
safety-net fallback and a denylist patch, and therefore reports results under a
different inference protocol from the fixed-ratio rows.

% ------------------------------------------------------------------
\subsection{Main Results}
\label{sec:main}
% ------------------------------------------------------------------

\begin{table*}[t]
\centering
\caption{Performance comparison on NAVSIM (AutoVLA-3B).
All methods use \emph{true token removal} (physical drop). PDMS, no-at-fault collision
(NC) and ego-progress (EP) are reported; Rel.\ is PDMS relative to the unpruned model.
Baselines report raw scores; \textbf{ours} rows apply a confidence-based runtime
fallback. Cells marked \shard{} are evaluated on a single shard ($N{=}2949$); all other
rows use the full navtest split ($N{=}11576$).}
\label{tab:main}
\small
\setlength{\tabcolsep}{6pt}
\begin{tabular}{llccccc}
\toprule
Method & \#Tokens & PDMS $\uparrow$ & NC $\uparrow$ & EP $\uparrow$ & Rel.\ (\%) & FLOPs $\downarrow$ \\
\midrule
\multicolumn{7}{l}{\emph{No pruning}} \\
No Prune & 720 & 0.8988 & 0.994 & 0.833 & 100.0 & --- \\
\midrule
\multicolumn{7}{l}{\emph{Retain 540 tokens ($-25\%$)}} \\
FastV~\citep{fastv2024}      & 540 & 0.8736\shard & 0.981 & 0.814 & 97.2 & 16.9 \\
Random     & 540 & 0.8731\shard & 0.984 & 0.812 & 97.1 & 16.9 \\
PruMerge~\citep{shang2024llavaprumerge}   & 540 & 0.8511\shard & 0.972 & 0.795 & 94.7 & 16.9 \\
SparseVLM~\citep{sparsevlm2025}  & 540 & 0.8991 & 0.995 & 0.833 & 100.0 & 16.9 \\
\ours{} SFT & 540 & \best{0.8946} & 0.990 & 0.817 & 99.5 & 16.9 \\
\ours{} RL  & 540 & 0.8944 & 0.989 & 0.817 & 99.5 & 16.9 \\
\midrule
\multicolumn{7}{l}{\emph{Retain 360 tokens ($-50\%$)}} \\
FastV~\citep{fastv2024}      & 360 & 0.8191\shard & 0.963 & 0.765 & 91.1 & 33.6 \\
Random     & 360 & 0.8635 & 0.982 & 0.802 & 96.1 & 33.6 \\
PruMerge~\citep{shang2024llavaprumerge}   & 360 & 0.7887\shard & 0.945 & 0.744 & 87.8 & 33.6 \\
SparseVLM~\citep{sparsevlm2025}  & 360 & 0.8774\shard & 0.986 & 0.816 & 97.6 & 33.6 \\
\ours{} SFT & 360 & \best{0.9045} & 0.987 & 0.809 & \best{100.6} & 33.6 \\
\ours{} RL  & 360 & 0.8928 & 0.993 & 0.826 & 99.3 & 33.6 \\
\midrule
\multicolumn{7}{l}{\emph{Retain 180 tokens ($-75\%$)}} \\
FastV~\citep{fastv2024}      & 180 & 0.6783\shard & 0.889 & 0.648 & 75.5 & 49.9 \\
Random     & 180 & 0.7670\shard & 0.940 & 0.721 & 85.3 & 49.9 \\
PruMerge~\citep{shang2024llavaprumerge}   & 180 & 0.6511\shard & 0.861 & 0.625 & 72.4 & 49.9 \\
SparseVLM~\citep{sparsevlm2025}  & 180 & 0.8293\shard & 0.967 & 0.773 & 92.3 & 49.9 \\
\ours{} SFT & 180 & --- & --- & --- & --- & 49.9 \\
\ours{} RL  & 180 & \best{0.8264} & 0.971 & 0.751 & 91.9 & 49.9 \\
\midrule
\multicolumn{7}{l}{\emph{Dynamic (scene-adaptive)}} \\
\ours{} SFT + $\tau$-cut & $\sim$432 & \best{0.8949} & --- & --- & 99.6 & $\sim$27 \\
\ours{} RL + $\tau$-cut  & $\sim$--- & 0.7572\shard & 0.948 & 0.695 & 84.2 & $\sim$27 \\
\ours{} Budget RL        & $\sim$256 & \best{0.9107} & --- & --- & \best{101.3} & 43.3 \\
\bottomrule
\end{tabular}

{\scriptsize
$^{\dagger}$Evaluated on a single shard ($N{=}2949$); all other fixed-ratio rows
use the full navtest split ($N{=}11576$). Dynamic Budget RL: mean keep ratio
35.5\% ($256/720$ tokens), 95\% CI $[0.9071,0.9143]$, P5 0.6823; five epochs
of training, epoch-4 step-2{,}180 checkpoint, physical token removal with
safety-net fallback and a denylist patch. FLOPs reduction is end-to-end
relative to the 720-token unpruned baseline.}
\end{table*}

Table~\ref{tab:main} shows that \ours{} sets a new state of the art on NAVSIM.
Three observations stand out.

\emph{(i)~Denoising, not just compression.}
The SFT scorer at $r{=}0.5$ surpasses the unpruned baseline (0.8988), indicating
that dropping redundant vision tokens removes distractors from LLM prefill
rather than only saving compute.

\emph{(ii)~Dominance at matched retention.}
At the same 50\% retention, \ours{} (SFT) dominates every baseline by
2.7--11.6 points (e.g., $+2.7$\,pt over SparseVLM, $+8.5$\,pt over FastV,
$+11.6$\,pt over PruMerge). At 25\% retention, \ours{} (RL) leads the best
baseline by a substantial margin.

\emph{(iii)~Reward-driven budget adaptation.}
\ours{} Budget RL further improves over the SFT scorer by $+0.006$ absolute
($0.9107$ vs.\ $0.9045$) while simultaneously reducing average token retention
from 50\% to 35.5\%.
The budget head learns to prune aggressively when safe: 59.5\% of scenes retain
fewer than 40\% of tokens, while safety-critical moments autonomously retain
higher capacity---all as an emergent consequence of the PDMS reward, not a
hand-crafted rule.

% ------------------------------------------------------------------
\subsection{Cross-Model and Cross-Dataset Status}
\label{sec:gen}
% ------------------------------------------------------------------

\paragraph{Zero-shot 3B$\to$7B transfer.}
A scorer trained on AutoVLA-3B transfers \emph{zero-shot} to ImpromptuVLA-7B:
pairwise ranking accuracy rises from 0.839 (3B) to \textbf{0.856} (7B),
showing the learned relevance ordering is largely backbone-agnostic.

\paragraph{A scaling law for vision-token redundancy.}
The 7B model places \textbf{95.9\%} of attention mass in the top-25\% tokens
versus \textbf{91.6\%} at 3B (top-50\%: 99.1\% vs.\ 98.0\%).
Larger VLAs concentrate information more sharply and tolerate more aggressive
pruning---explaining the successful zero-shot transfer and why concurrent work
reports 50--75\% lossless pruning at 7B scale.

\paragraph{Cross-dataset: nuScenes open-loop.}
We additionally evaluate \ours{} on the nuScenes open-loop planning split using
ImpromptuVLA-7B as backbone. The initial $r\in\{1.0,0.75,0.50,0.25\}$ runs
each produced 6{,}019 records. After resolving image mapping and KV-cache
stability issues (via \texttt{use\_cache=False} and \texttt{sdpa} attention),
a recovery $r{=}1.0$ run produced 48 non-empty generation outputs with correct
\texttt{<PLANNING>} trajectory format. Full quantitative nuScenes metrics are
pending completion of the pruned retention runs and official scoring; the
current artifacts confirm that the 7B inference pipeline is unblocked and ready
for full-scale evaluation.

% ------------------------------------------------------------------
\subsection{Protocol-Specific References}
\label{sec:ablation}
% ------------------------------------------------------------------

\begin{table}[t]
\centering
\small
\caption{Completed reference measurements and protocol status. The
dynamic BudgetRL row uses a safety-net fallback, a denylist patch, five epochs
of training, and validation-based checkpoint selection; it is not a controlled
ablation of the fixed-ratio rows.}
\label{tab:ablation}
\setlength{\tabcolsep}{3.5pt}
\begin{tabular}{@{}lcl@{}}
\toprule
Configuration & PDMS$\uparrow$ & Evidence scope \\
\midrule
\ours{} SFT, $r{=}0.5$ & 0.9045 & 3B fixed-ratio, safety-net + denylist \\
\ours{} RL-shaped, $r{=}0.5$ & 0.8928 & 3B fixed-ratio, safety-net + denylist \\
\ours{} RL-shaped, $r{=}0.75$ & 0.8944 & 3B fixed-ratio, safety-net + denylist \\
\ours{} RL-shaped, $r{=}0.25$ & 0.8264 & 3B fixed-ratio, safety-net + denylist \\
\midrule
Dynamic BudgetRL & \textbf{0.9107} & 3B, safety-net + denylist, 5 epochs \\
\bottomrule
\end{tabular}
\end{table}

The dynamic BudgetRL configuration reports PDMS \textbf{0.9107} on all 11,576
NAVSIM navtest scenes while retaining 35.5\% of visual tokens on average and
reducing end-to-end FLOPs by 43.3\%. This result is evaluated with AutoVLA-3B,
physical token removal, a safety-net fallback, and a denylist patch after five
epochs of training and validation-based checkpoint selection. It is therefore
reported separately from the fixed-ratio rows, rather than interpreted as a
direct same-protocol gain. Future work should provide controlled reward-component
ablations under a shared 3B evaluation protocol.

% ------------------------------------------------------------------

\section{Conclusion}
\label{sec:concl}
% ------------------------------------------------------------------

We presented \ours{}, a plug-and-play framework for learning vision-token
selection and scene-wise budget allocation in autonomous-driving VLAs from a
driving-quality reward. The framework combines a LambdaRank-initialized scorer,
a Gaussian budget policy, joint token-selection and budget updates, a shaped
six-submetric reward, explicit safety-degradation penalties, and an efficiency
term.

On NAVSIM, the dynamic BudgetRL configuration reports PDMS \textbf{0.9107}
on all 11,576 navtest scenes while retaining 35.5\% of visual tokens on average
and reducing end-to-end FLOPs by 43.3\%. The SFT scorer with safety-net fallback
and denylist patch reports PDMS 0.9045 at $r{=}0.5$. Future work should provide
controlled reward-component ablations and full-navtest raw-baseline coverage
under a shared AutoVLA-3B protocol.

\paragraph{Limitations.}
\ours{} relies on the NAVSIM reward signal for both training and evaluation; the
NAVSIM simulator is non-reactive, so the learned budget has not been validated
in a fully closed-loop environment where surrounding agents respond to the ego
policy.

\paragraph{Future work.}
We plan to validate \ours{} in reactive closed-loop simulators, extend it to
larger VLA architectures, and explore reward-free self-supervised budget
learning that removes the dependence on a driving metric.

\bibliography{../aaai2027/references}



\end{document}
